\documentclass[11pt,letterpaper]{article}
\usepackage[T1]{fontenc}
\usepackage[utf8]{inputenc}
\usepackage[margin=0.88in,top=0.84in,bottom=0.83in,headheight=13pt,headsep=20pt,footskip=28pt]{geometry}
\usepackage{newpxtext}
\usepackage{amsmath,amsthm}
\usepackage{newpxmath}
\usepackage{microtype}
\usepackage{xcolor}
\definecolor{Forest}{HTML}{30392C}
\definecolor{Olive}{HTML}{687144}
\definecolor{Muted}{HTML}{65685F}
\definecolor{Sage}{HTML}{CBD0BC}
\definecolor{Pale}{HTML}{F2F4EB}
\usepackage{mathtools}
\usepackage{booktabs,tabularx,longtable,array}
\usepackage{enumitem}
\usepackage{titlesec}
\usepackage{fancyhdr}
\usepackage[most]{tcolorbox}
\usepackage{needspace}
\PassOptionsToPackage{obeyspaces}{url}
\usepackage{xurl}
\usepackage[colorlinks=true,linkcolor=Olive,citecolor=Olive,urlcolor=Olive,bookmarksnumbered=true,pdfencoding=auto,psdextra]{hyperref}
\hypersetup{pdftitle={Finite definable choice from one measurable cardinal},pdfauthor={Prepared for Vladimir Reshetnikov},pdfsubject={Prikry tail recoding, finite equivariant rules, and a sharp local ultrafilter threshold},pdfkeywords={Prikry forcing, measurable cardinal, definability, finite choice, ultrafilters, HOD}}
\urlstyle{same}
\setlength{\parindent}{0pt}
\setlength{\parskip}{5.1pt plus 1pt minus .4pt}
\linespread{1.035}
\setlength{\emergencystretch}{2em}
\setcounter{tocdepth}{1}
\setcounter{secnumdepth}{2}
\titleformat{\section}{\Large\bfseries\color{Forest}}{\thesection}{.6em}{}
\titleformat{\subsection}{\large\bfseries\color{Forest}}{\thesubsection}{.55em}{}
\titleformat{\subsubsection}{\normalsize\bfseries\color{Forest}}{}{0pt}{}
\titlespacing*{\section}{0pt}{18pt plus 3pt minus 2pt}{8pt}
\titlespacing*{\subsection}{0pt}{12pt plus 2pt minus 1pt}{5pt}
\titlespacing*{\subsubsection}{0pt}{9pt}{3pt}
\setlist[itemize]{leftmargin=1.4em,itemsep=2pt,topsep=3pt,parsep=0pt}
\setlist[enumerate]{leftmargin=1.7em,itemsep=3pt,topsep=4pt,parsep=0pt}
\pagestyle{fancy}
\fancyhf{}
\fancyhead[L]{\sffamily\fontsize{9}{11}\selectfont\color{Muted}LARGE CARDINALS AT THE INCONSISTENCY FRONTIER}
\fancyhead[R]{\sffamily\fontsize{9}{11}\selectfont\color{Muted}CONTINUATION\quad 18 SEP 2026}
\fancyfoot[L]{\small\color{Muted}Finite symmetry from one measurable cardinal}
\fancyfoot[R]{\small\color{Muted}\thepage}
\renewcommand{\headrulewidth}{.35pt}
\renewcommand{\footrulewidth}{0pt}
\fancypagestyle{plain}{\fancyhf{}\fancyfoot[L]{\small\color{Muted}Finite symmetry from one measurable cardinal}\fancyfoot[R]{\small\color{Muted}\thepage}\renewcommand{\headrulewidth}{0pt}}
\tcbset{colback=Pale,colframe=Sage,colbacktitle=Sage,coltitle=Forest,fonttitle=\bfseries,boxrule=.5pt,arc=3pt,left=9pt,right=9pt,top=6pt,bottom=6pt,toptitle=3pt,bottomtitle=3pt,before skip=8pt,after skip=8pt,breakable}
\newtcolorbox{assessment}[1]{title={#1}}
\theoremstyle{definition}
\newtheorem{definition}{Definition}[section]
\newtheorem{theorem}[definition]{Theorem}
\newtheorem{proposition}[definition]{Proposition}
\newtheorem{lemma}[definition]{Lemma}
\newtheorem{corollary}[definition]{Corollary}
\newtheorem{remark}[definition]{Remark}
\newtheorem{question}[definition]{Question}
\newtheorem{inputthm}{Input}
\renewcommand{\theinputthm}{\Alph{inputthm}}
\numberwithin{equation}{section}
\newcommand{\ZFC}{\mathsf{ZFC}}
\newcommand{\ZF}{\mathsf{ZF}}
\newcommand{\AC}{\mathsf{AC}}
\newcommand{\GA}{\mathsf{GA}}
\newcommand{\HOD}{\mathrm{HOD}}
\newcommand{\OD}{\mathrm{OD}}
\newcommand{\CD}{\mathrm{CD}}
\newcommand{\HCD}{\mathrm{HCD}}
\newcommand{\UF}{\mathrm{UF}}
\newcommand{\Ex}{\operatorname{Ex}}
\newcommand{\UEx}{\operatorname{UEx}}
\newcommand{\Ext}{\operatorname{Ext}}
\newcommand{\SC}{\operatorname{SC}}
\newcommand{\CEx}{\operatorname{CEx}}
\newcommand{\REx}{\operatorname{REx}}
\newcommand{\Con}{\operatorname{Con}}
\newcommand{\crit}{\operatorname{crit}}
\newcommand{\cf}{\operatorname{cf}}
\newcommand{\ot}{\operatorname{ot}}
\newcommand{\ran}{\operatorname{ran}}
\newcommand{\rank}{\operatorname{rank}}
\newcommand{\lcm}{\operatorname{lcm}}
\newcommand{\Ord}{\mathrm{Ord}}
\newcommand{\Pow}{\mathcal P}
\newcommand{\id}{\mathrm{id}}
\newcommand{\restr}{\mathbin{\upharpoonright}}
\newcommand{\Izero}{\mathrm{I0}}
\newcommand{\Itwo}{\mathrm{I2}}
\newcommand{\Ithree}{\mathrm{I3}}
\newcommand{\Iwf}{\mathrm{I3}_{\mathrm{wf}}(0)}
\newcommand{\Sel}{\operatorname{Sel}}
\newcommand{\FF}{\mathcal F}
\newcommand{\PP}{\mathbb P}
\newcommand{\Clam}{\mathcal C_\lambda}
\newcommand{\Rig}{\operatorname{Rig}}
\newcommand{\Spec}{\operatorname{Spec}}
\newcommand{\ch}{\operatorname{ind}}
\newcommand{\CF}{\mathsf{CF}}
\newcommand{\src}[1]{\unskip\ {\normalfont\small\sffamily\color{Olive}[#1]}\ }
\newcommand{\R}[1]{\textsf{R#1}}
\newcolumntype{Y}{>{\raggedright\arraybackslash}X}
\newcolumntype{P}[1]{>{\raggedright\arraybackslash}p{#1}}
\newcolumntype{C}{>{\centering\arraybackslash}p{.034\textwidth}}
\newcommand{\yes}{$\bullet$}
\newcommand{\prt}{$\circ$}
\newcommand{\arxiv}[1]{\href{https://arxiv.org/abs/#1}{arXiv:#1}}
\newcommand{\doi}[1]{\href{https://doi.org/#1}{doi:\nolinkurl{#1}}}


\newcommand{\Dk}{D_\kappa}
\newcommand{\Qk}{Q_\kappa}
\newcommand{\sq}{q_*}
\newcommand{\Gd}{G_d}
\newcommand{\Fk}{\operatorname{Fix}}
\newcommand{\Tail}{\operatorname{Tail}}
\newcommand{\Sym}{\operatorname{Sym}}
\newcommand{\fin}{\mathrm{fin}}
\newcommand{\eps}{\varepsilon}
\newcommand{\ttuple}[2]{{#1}^{\langle #2\rangle}}
\newcommand{\wordeq}{\mathrel{\sim_{\mathrm{st}}}}
\newcommand{\tr}{\operatorname{tr}}
\newcommand{\seq}[1]{\langle #1\rangle}
\newcommand{\Eone}{\mathsf{FS}}
\newcommand{\sh}{\operatorname{sh}}
\begin{document}
\thispagestyle{empty}
{\sffamily\bfseries\small\color{Olive}SET THEORY\quad /\quad FOURTH CONTINUATION}

\vspace{13pt}
{\fontsize{28}{31}\selectfont\bfseries\color{Forest}Finite definable choice\\from one measurable cardinal\par}
\vspace{10pt}
{\fontsize{14}{18}\selectfont Prikry tail recoding, finite equivariant rules,\\and a sharp local ultrafilter threshold\par}
\vspace{14pt}
{\color{Muted}Prepared for Vladimir Reshetnikov\hfill 18 September 2026\par}
\vspace{3pt}
{\color{Sage}\hrule height .6pt}
\vspace{12pt}

\textbf{Abstract.}
Let $W$ satisfy $\ZFC$, let $U$ be a normal measure on $\kappa$ in $W$, and let
$V=W[c]$ be its Prikry extension. Write $q_*=[c]$ for the class of the generic
cofinal sequence modulo finite symmetric difference. We prove that, even allowing
$q_*$ and \emph{arbitrary ground-model set parameters}, there is no nonempty
finite definable family of proper selectors on finite subsets of
$Q_\kappa=D_\kappa/{=^*}$. More generally, a finite permutation-equivariant
labeling problem has a nonempty finite definable family of solutions if and
only if each individual permutation fixes a label; when this condition holds,
a single solution is definable from a low-rank parameter.

This answers the Prikry-extension part of the explicit third-round question
in the supplied synthesis. The proof replaces ultraexacting embeddings by a
same-universe comparison: deleting one sufficiently late generic point
preserves the generic tail class and the entire extension, while rotating
finite residue configurations. Deciding the \emph{set of finite tables},
rather than a chosen member of a finite family, makes the argument work
without an illicit definable choice.

We also prove that definable traces on $q_*$ are clopen in its relative Prikry
topology and obtain full cardinality in every integer phase, with full-sized
coordinate projections. Finally, no nonempty finite definable family of
ultrafilters on $q_*$ exists, but a countably infinite such family does exist.
The latter construction is a general $\ZFC$ fact about tail classes. It is
not a construction of countably many global ultrafilter-valued kernels.

\begin{assessment}{Main advance and limits}
The finite-selector and finite-label \emph{obstructions}, previously obtained
in the archive from ultraexactingness, already occur in an ordinary Prikry
extension of one measurable cardinal. Their proofs even tolerate the
parameter $q_*$ and every fixed ground-model parameter.

This is a result about a specific Prikry model, not a proof that abstract
rigidity alone implies these obstructions. No inconsistency of an ordinary
large-cardinal axiom, no new lower consistency bound, and no worldwide
priority claim are asserted.
\end{assessment}

\vfill
{\small\color{Muted}Detailed proofs in English. The accompanying finite tests check only
finite permutation algebra; the new set-theoretic arguments are not Lean formalizations.}
\newpage
\tableofcontents
\newpage

\section{The result and its place in the research program}\label{sec:overview}

\subsection{The question that is answered}

The supplied \emph{Large Cardinals Synthesis} \cite{archive} distinguishes two
sources of definability obstructions. Its rigidity arguments already have a
one-measurable Prikry construction. Its finite-family selector theorem and
the necessity direction of the finite-label classification were proved using
ultraexacting embeddings. The last question of its Section~15 asks whether
these latter conclusions also hold in a Prikry extension of a measurable.
The answer proved here is affirmative.

The source question also describes this as asking whether the conclusions
follow from rigidity alone. Those are different logical assertions. A theorem
in every normal Prikry extension establishes the first, but does not establish
the second. Our proofs use a concrete finite-tail recoding property, not just
regularity in a relativized $\HOD$.

The distinction matters for interpreting the large-cardinal frontier.
Ultraexactingness has the consistency strength of $\Izero$ by
Aguilera--Bagaria--Goldberg--L\"ucke \cite{abgl}. The present theorem shows that
these particular finite-choice consequences have a construction using only
one measurable cardinal. It does not reduce the consistency strength of
ultraexactingness itself.

\subsection{Statement in quotient language}

For an uncountable cardinal $\lambda$ of cofinality $\omega$, put
\[
 D_\lambda=\{a\subseteq\lambda:\ot(a)=\omega,
                          \ \sup a=\lambda\},
 \qquad Q_\lambda=D_\lambda/{=^*},
\]
where $a=^*b$ means that $a\triangle b$ is finite. We identify an element of
$D_\lambda$ with its increasing enumeration when writing $a(n)$. For a set
$X$, the notation $\ttuple{X}{m}$ means the ordered $m$-tuples of pairwise
distinct elements of $X$.

A parameter $q\in Q_\lambda$ means the equivalence class as \emph{one set},
not a chosen representative and not permission to use all its members as
parameters. Thus $\OD_{p,q}$ denotes definability from $p,q$, and finitely many
ordinals. For $\OD_{V_\lambda,q}$, finitely many individual parameters in
$V_\lambda$ are allowed. All definability and cardinalities below are computed
in the final universe $V$, unless another model is displayed.

Finite labeling data are coded by hereditarily finite sets: we take the
finite label set $F$ to be a finite subset of $\omega$ with the specified
action. This convention is part of the low-rank definability assertions;
renaming arbitrary labels would also require their names as parameters.

\begin{theorem}[Main theorem]\label{thm:main}
Let $W\models\ZFC$, let $U\in W$ be a normal measure on $\kappa$, and let
$V=W[c]$ be the normal Prikry extension. Let $q_*=[c]\in Q_\kappa$.
Then the following assertions hold in $V$.
\begin{enumerate}
\item For every $p\in W$, every $n\ge2$, and every $0<r<n$, there is no
nonempty finite $\OD_{p,q_*}$ family of functions
\[
 f:[Q_\kappa]^n\longrightarrow[Q_\kappa]^r,
 \qquad f(B)\subseteq B.
\]
\item Let $m\ge1$, let $\Gamma\le S_m$, and let $F$ be a finite
$\Gamma$-set. The following are equivalent:
\begin{enumerate}[label=(\alph*)]
\item for some $p\in W$, there is a nonempty finite $\OD_{p,q_*}$ family of
$\Gamma$-equivariant maps $\ttuple{Q_\kappa}{m}\to F$;
\item there is one such map in $\OD_{V_\kappa}$;
\item $\Fk_F(g)\ne\varnothing$ for every $g\in\Gamma$.
\end{enumerate}
\item Every $\OD_{p,q_*}$ subset of $q_*$, for $p\in W$, is clopen in the
relative Prikry topology. If it is nonempty, then it has $\kappa$ members in
every integer phase relative to any $a\in q_*$. Eventually, each coordinate
projection of each phase has cardinality $\kappa$.
\item The least size of a nonempty $\OD_{V_\kappa,q_*}$ family of ultrafilters
on the \emph{full} power set of $q_*$ is exactly $\aleph_0$.
\end{enumerate}
Moreover $\kappa$ is a strong limit of cofinality $\omega$ in $V$,
$V_\kappa^V=V_\kappa^W$, and $q_*$ is rigid in the sense of the supplied
synthesis.
\end{theorem}

Part~(2) concerns quotient tuples. A function on tuples of representatives
that is invariant under finite changes of each coordinate factors uniquely
through these quotient tuples. Thus this is the same labeling problem as in
the source reports. The extension here is both the Prikry-model proof of
necessity and the allowance of a \emph{finite family} of candidate rules.

\begin{center}
\small
\renewcommand{\arraystretch}{1.15}
\begin{tabularx}{\textwidth}{@{}P{.26\textwidth}YY@{}}
\toprule
\textbf{Topic}&\textbf{Supplied synthesis}&\textbf{This continuation}\\\midrule
Finite selectors&Ultraexacting proof; Prikry case explicitly open&No finite
family in every normal Prikry extension, even relative to $q_*$ and ground parameters\\
Finite labels&Single-rule classification at an ultraexacting cardinal&Same
criterion in a Prikry extension, also for finite families\\
Full fibres&Embedding-based phase and coordinate analysis&Relative clopenness
and full-sized phase-coordinate projections at $q_*$\\
Ultrafilter selection&Finite obstructions; countable global-kernel question&Exact
countable threshold for \emph{local families}, with the global question kept separate\\
\bottomrule
\end{tabularx}
\end{center}

\section{Prikry forcing and the same-universe recoding lemma}\label{sec:recoding}

\subsection{Classical input and conventions}

Here a normal measure is nonprincipal and $\kappa$-complete, with the usual
normality property for regressive functions. A condition in $\PP_U^W$ is
$(s,A)$, where $s$ is a finite increasing sequence
of ordinals below $\kappa$, $A\in U$, and all members of $A$ exceed the last
member of $s$. The empty stem is permitted. We use the convention that
$(t,B)\le(s,A)$ is stronger: $t$ end-extends $s$, its new entries belong to $A$,
and $B\subseteq A$. For an increasing cofinal sequence $d$, write
\[
 G_d=\{(s,A)\in\PP_U^W:
            s\text{ is an initial segment of }d,
            \ \ran(d)\setminus\ran(s)\subseteq A\}.
\]
We call this the filter \emph{realized by} $d$.

\begin{inputthm}[Normal Prikry forcing]\label{in:prikry}
Normal Prikry forcing preserves cardinals, changes $\cf(\kappa)$ to $\omega$,
and adds no bounded subsets of $\kappa$. For an increasing cofinal
$d:\omega\to\kappa$, the Mathias criterion says
\[
 G_d\text{ is generic over }W
 \quad\Longleftrightarrow\quad
 \text{for every }A\in U,
       \ d(n)\in A\text{ for all sufficiently large }n.
\]
The extension is $W[G_d]=W[d]$.
\end{inputthm}

The criterion is due to Mathias \cite{mathias}. A convenient source for the
preservation statements and the criterion is Benhamou
\cite[Corollaries~3.12, 3.15; Theorem~3.17]{benhamou}; the numbering here refers
to the cited arXiv version. These are the only specialized forcing facts
imported in the new proofs. We also use the ordinary forcing theorem and
truth lemma.

Since a measurable cardinal is strongly inaccessible in its ground model,
these facts give $V_\kappa^W=V_\kappa^V$. Here is the small extra argument,
which is sometimes obscured by the phrase ``no bounded subsets.'' Induct on
$\alpha<\kappa$. If $V_\alpha^W=V_\alpha^V$, a ground-model bijection of
$V_\alpha^W$ with an ordinal below $\kappa$ codes every subset of that rank by
a bounded subset of $\kappa$. Hence no new subset of $V_\alpha$ appears. At
limit ranks take unions. The required ground-model size bounds follow from
inaccessibility. The same preservation facts show that $\kappa$ remains a
strong limit cardinal.

\subsection{Finite changes preserve the whole extension}

\begin{lemma}[Finite-tail recoding]\label{lem:recoding}
Suppose $d\in D_\kappa$ and $d=^*c$.
Then $d$ is Prikry generic over $W$,
\[
                 W[d]=W[c]=V,\qquad [d]=[c]=q_*.
\]
If $P=(s,A)\in G_c$ and $d$ is obtained by deleting a point of $c$ after its
stem $s$, then $P\in G_d$ as well.
\end{lemma}

\begin{proof}
For every $A\in U$, the sequence $c$ is eventually in $A$. A finite change
preserves this assertion, so the Mathias criterion applies to $d$. The
finite sets $c\setminus d$ and $d\setminus c$ consist of ordinals and therefore
belong to the transitive ground model $W$. Consequently $c\in W[d]$ and
$d\in W[c]$, giving equality of the two extensions. The quotient classes
are equal by their definition in this common universe.

For the last assertion, deleting a point beyond the stem leaves $s$ an
initial segment and leaves every remaining point beyond $s$ in $A$.
This is exactly the condition $P\in G_d$.
\end{proof}

\begin{remark}[What the lemma does not say]\label{rem:noautomorphism}
There is no asserted nonidentity automorphism of the transitive universe
$V$, and no elementary embedding of $V$ is being constructed. We compare two
generic filters for the same ground-model forcing that produce the
\emph{same} extension. A set defined in that extension from fixed ground
parameters, ordinals, and $q_*$ is literally the same set under both
presentations. The generic sequence used to evaluate that set can change.
\end{remark}

\begin{lemma}[Deciding a finite observation]\label{lem:decision}
Fix ground parameters $p$ and ordinals $\vec\alpha$. Suppose a formula
uniquely defines an object $Z$ in $V$ from $p,q_*,\vec\alpha$. Suppose a
finite observation $\mathcal O(d,Z)$ is defined from the increasing sequence
$d$ and $Z$, and always belongs to a fixed finite ground-model set $E$ in a
suitable forcing neighborhood of $G_c$.
Then there are $P=(s,A)\in G_c$ and $e\in E$ such that
\[
              \mathcal O(d,Z)=e
\]
for every $d=^*c$ with $P\in G_d$. Here the object $Z$ is the same actual
object in $V$ throughout.
\end{lemma}

\begin{proof}
By the truth lemma, choose a condition in $G_c$ forcing the stated uniqueness
and the properties needed to form the observation. Replace $q_*$ by the
canonical name $[\dot c]$ in its definition. Strengthen inside $G_c$ to a
condition $P$ deciding the finite-valued name for $\mathcal O(\dot c,Z)$,
say with value $e$.

If $d=^*c$ realizes $P$, Lemma~\ref{lem:recoding} says that $d$ is generic
and that its extension and tail-class parameter are the original $V,q_*$.
The forced uniqueness thus identifies the object defined using $[d]$ with
the original $Z$. Since $P\in G_d$, the decided observation has the value
$e$ there too.
\end{proof}

This lemma is the replacement for the embedding interface in the source
reports. The observation will be a \emph{set of finite tables}. In particular,
we will not choose a definable member of a definable finite family. Such a
choice is precisely what must not be assumed.

\section{Finite cyclic configurations and finite families}\label{sec:tables}

\subsection{Realizing an arbitrary finite permutation}

Let $\Omega$ be a finite set and $\tau$ a permutation of $\Omega$. Enumerate
its cycles as
\[
 C_j=(u_{j,0},\ldots,u_{j,d_j-1}),\qquad j<b,
 \qquad \tau(u_{j,i})=u_{j,i+1\bmod d_j}.
\]
For an increasing cofinal sequence $a$, define
\begin{equation}\label{eq:cyclecode}
 A_{j,i}(a)=\{a(d_j k+i)+j:k<\omega\},\qquad
 x_{u_{j,i}}(a)=[A_{j,i}(a)]\in Q_\kappa.
\end{equation}
The addition is ordinal addition by a finite number. It remains below the
uncountable cardinal $\kappa$.

\begin{lemma}[Residues with distinct offsets]\label{lem:cycles}
For $a=^*c$, the classes $x_u(a)$, $u\in\Omega$, are pairwise distinct.
If $d$ is obtained from $a$ by deleting one point, then
\begin{equation}\label{eq:rotation}
                        x_u(d)=x_{\tau u}(a)
                        \qquad(u\in\Omega).
\end{equation}
Both statements are unaffected by any exceptional finite initial segment.
\end{lemma}

\begin{proof}
A normal measure contains the club of infinite limit ordinals, so all
sufficiently late entries of $a$ are limit ordinals. The finitely many
translations by $j<b$ are then separated: if $\alpha<\beta$ are two such
limit ordinals, $\alpha+j<\beta$ for every $j<b$. For a fixed entry of $a$,
different offsets give different ordinals. Hence the sets assigned to
different cycles are eventually disjoint. Within a cycle, different
residue classes use disjoint entries of $a$. All the sets in
\eqref{eq:cyclecode} are cofinal of order type $\omega$. Distinct sets in
this construction therefore have infinite symmetric difference.

If the deleted entry has index $h$, then $d(n)=a(n+1)$ for all $n\ge h$.
For a residue $i<d_j-1$, this advances the residue to $i+1$. For the last
residue, it advances to residue $0$ in the next block. The discarded initial
pieces are finite. Taking quotient classes gives \eqref{eq:rotation}.
\end{proof}

The normal-measure assertion used here has an elementary verification. If a
club $C$ were not measure one, the map
$\alpha\mapsto\sup(C\cap\alpha)$ would be regressive on the measure-one set
of sufficiently large $\alpha\notin C$. Normality would make it constant on
a measure-one set, which would then be bounded below the next point of $C$.
This contradicts nonprincipality and $\kappa$-completeness.

\subsection{A family is represented by a set of tables}

Let $m\ge1$ and let $F$ be finite. Suppose $\mathscr L$ is a nonempty finite
$\OD_{p,q_*}$ family of maps
\[
                         L:\ttuple{Q_\kappa}{m}\longrightarrow F,
                         \qquad p\in W.
\]
For the configuration from \eqref{eq:cyclecode}, define
\[
 t_L^a:\ttuple{\Omega}{m}\longrightarrow F,
 \qquad
 t_L^a(u_0,\ldots,u_{m-1})
     =L(x_{u_0}(a),\ldots,x_{u_{m-1}}(a)),
\]
and let
\[
                       T(a)=\{t_L^a:L\in\mathscr L\}.
\]
Distinct maps can have the same finite table. This causes no problem:
$1\le |T(a)|\le |\mathscr L|$.

\begin{lemma}[Finite-table symmetry]\label{lem:table}
Let $M=|\mathscr L|$ and let $\Omega,\tau$ be any finite configuration with
$|\Omega|\ge m$. The set $T(c)$ of tables is invariant under the permutation
\[
 R_\tau(t)(u_0,\ldots,u_{m-1})
                  =t(\tau u_0,\ldots,\tau u_{m-1}).
\]
Consequently, for $H=\lcm(1,2,\ldots,M)$,
\begin{equation}\label{eq:tablefixed}
 t(\tau^H u_0,\ldots,\tau^H u_{m-1})=t(u_0,\ldots,u_{m-1})
 \quad\text{for every }t\in T(c).
\end{equation}
\end{lemma}

\begin{proof}
The set of all maps $\ttuple{\Omega}{m}\to F$ is finite. Its power set is
also a finite ground-model set. Apply Lemma~\ref{lem:decision} to the
observation $T(a)$, with $Z=\mathscr L$. Obtain a condition $P\in G_c$
deciding $T(\dot c)$ to be the actual set $T(c)$.

Delete one entry of $c$ after the stem of $P$, obtaining $d$. Then $d$
realizes $P$, so $T(d)=T(c)$. On the other hand, the family $\mathscr L$ is
unchanged, while Lemma~\ref{lem:cycles} gives
\[
                          t_L^d=R_\tau(t_L^c)
                          \qquad(L\in\mathscr L).
\]
Thus $T(c)=R_\tau[T(c)]$. The map $R_\tau$ is injective on all tables, so its
restriction to the nonempty finite set $T(c)$ is a permutation.

Every cycle of this permutation has length at most $|T(c)|\le M$ and hence
a length dividing $H$. Therefore $R_\tau^H$ fixes each individual table.
The explicit formula for this assertion is \eqref{eq:tablefixed}.
\end{proof}

\begin{remark}[Why the exponent is indispensable]
A finite definable family can be permuted without any one of its members
being fixed. Already two binary choices on a two-point set are exchanged by
the transposition. It is invalid to deduce a contradiction merely from a
nontrivial cycle on the inputs. The exponent $H$ kills the permutation of
the finite \emph{set of tables}; we then arrange that $\tau^H$ still acts
nontrivially on a suitable input. This is the finite-family amplification
principle of the source archive, now supplied with a forcing interface.
\end{remark}

\subsection{Amplification for any prescribed permutation}

Use the standard left action of $\Gamma\le S_m$ on ordered tuples:
\[
               (g\cdot\vec z)_i=z_{g^{-1}(i)}.
\]
A rule is equivariant when $L(g\cdot\vec z)=g\cdot L(\vec z)$.

\begin{theorem}[Elementwise fixed labels are necessary, even for finite families]
\label{thm:necessity}
If a nonempty finite $\OD_{p,q_*}$ family of equivariant maps
$\ttuple{Q_\kappa}{m}\to F$ exists, for $p\in W$, then every $g\in\Gamma$
has a fixed point in $F$.
\end{theorem}

\begin{proof}
Fix $g\in\Gamma$. Write its cycles on $\{0,\ldots,m-1\}$ as
\[
 (i_{j,0},\ldots,i_{j,\ell_j-1}),\qquad j<b,
 \quad g(i_{j,s})=i_{j,s+1\bmod\ell_j}.
\]
Let $M$ be the size of the family and $H=\lcm(1,\ldots,M)$. Construct a
finite set $\Omega$ with cycles
\[
 C_j=\{(j,t):t\in\mathbb Z/(H\ell_j)\mathbb Z\},
 \qquad\tau(j,t)=(j,t+1).
\]
Define an injection $v:m\to\Omega$ by
\[
                         v(i_{j,s})=(j,Hs).
\]
It has the crucial property
\begin{equation}\label{eq:amp}
                         \tau^H v(i)=v(g(i)).
\end{equation}
Realize $\Omega,\tau$ using Lemma~\ref{lem:cycles}, and use
Lemma~\ref{lem:table} for the given family. Choose any one table $t\in T(c)$;
this is an ordinary existential choice in the proof, not a definable
selection from the family. Let
\[
                           y=t(v(0),\ldots,v(m-1))\in F.
\]
The table inherits coordinate equivariance from its underlying rule. By
\eqref{eq:tablefixed} and \eqref{eq:amp},
\begin{align*}
 y
 &=t(\tau^H v(0),\ldots,\tau^H v(m-1))\\
 &=t(v(g(0)),\ldots,v(g(m-1)))\\
 &=g^{-1}\cdot t(v(0),\ldots,v(m-1))
   =g^{-1}\cdot y.
\end{align*}
The inverse in the third line comes from our displayed convention for the
left action on tuples. This equality is equivalent to $g\cdot y=y$.
Since $g$ was arbitrary, the conclusion follows.
\end{proof}

The forcing decision is made only after the finite amplified configuration
has been selected. Thus its stem length, however large, cannot spoil the
argument: one deletes a point still later. No uniform bound on deciding
conditions is needed.

\section{The exact finite-label classification}\label{sec:classification}

The sufficiency argument is present in the supplied synthesis
\cite[Section~12.5]{archive}. We include its full proof to make the new
Prikry classification self-contained and to delimit the new contribution.
Its use of cyclic stabilizers also explains why ``every element has a fixed
point'' cannot be replaced by ``the whole group has a common fixed point.''

\subsection{Incidence words and their stabilizers}

Let $\vec a=(a_0,\ldots,a_{m-1})$ be representatives of pairwise different
classes in $Q_\lambda$. The finite union $\bigcup_i a_i$ has order type
$\omega$: below each ordinal smaller than $\lambda$, each $a_i$ has only
finitely many points. Enumerate the union increasingly as
$\seq{\delta_n:n<\omega}$, and form the word
\[
                  w_{\vec a}(n)=\{i<m:\delta_n\in a_i\}
                  \in\Pow(m)\setminus\{\varnothing\}.
\]
Each label occurs infinitely often, and each pair of distinct labels is
separated infinitely often. The latter assertion expresses the fact that
$a_i\triangle a_j$ is infinite.

Call words $w,z$ \emph{shift-tail equivalent}, written $w\wordeq z$, if
there exist $r,s<\omega$ such that
$w(r+n)=z(s+n)$ for every $n<\omega$. Finite changes in the representatives
change the incidence word only by this equivalence: past the largest edited
ordinal, the unions and their incidence labels agree. Permuting coordinates
commutes with taking the incidence word.

\begin{lemma}[Cyclic stabilizers of word classes]\label{lem:stabilizer}
Let $w$ be a word with the two infinite-occurrence and separation properties
just stated. The stabilizer of $[w]_{\mathrm{st}}$ under any finite
$\Gamma\le S_m$ is cyclic.
\end{lemma}

\begin{proof}
Consider the subgroup
\[
 B_w=\{(d,g)\in\mathbb Z\times\Gamma:
             w(n+d)=g\cdot w(n)\text{ for all sufficiently large }n\}.
\]
The requirement ``sufficiently large'' ensures that negative shifts are
well-defined. Composition and inversion of eventual identities show
explicitly that $B_w$ is a subgroup of the direct product.

Its projection to $\mathbb Z$ has trivial kernel. Indeed, if $(0,g)\in B_w$,
then eventually $w(n)=g\cdot w(n)$. If $g$ moved $i$, the labels $i$ and
$g^{-1}(i)$ would have the same membership in every sufficiently late
$w(n)$, contradicting infinite separation. Hence $g$ is the identity.

The projection identifies $B_w$ with a subgroup of $\mathbb Z$. Such a
subgroup is cyclic, including the possible trivial subgroup. Its image
under the second projection is exactly the stabilizer of
$[w]_{\mathrm{st}}$: an eventual identity with an integer shift is equivalent
to a shift-tail equivalence. An image of a cyclic group is cyclic.
\end{proof}

\subsection{Constructing a rule when the test passes}

\begin{proposition}[Sufficiency in $\ZFC$]\label{prop:sufficiency}
Let $\lambda$ be any uncountable cardinal of cofinality $\omega$. Let
$\Gamma\le S_m$ act on a finite set $F$, and suppose every $g\in\Gamma$
fixes some element of $F$. Then an equivariant rule
$\ttuple{Q_\lambda}{m}\to F$ is definable from a well-order of the reals and
finite codes for $\Gamma,F$. These parameters belong to $V_\lambda$.
\end{proposition}

\begin{proof}
Words over the finite alphabet $\Pow(m)\setminus\{\varnothing\}$ are coded
by reals in a fixed way. Fix a well-order $\prec$ of the reals and a finite
ordering of $F$. All these objects have rank below $\omega+\omega$, hence
below the uncountable cardinal $\lambda$.

For each $\Gamma$-orbit of shift-tail classes of valid words, choose the
class $z_0$ containing the $\prec$-least real code in the union of those
classes. This is a definable choice from the real codes, not a choice of
cofinal representatives. By Lemma~\ref{lem:stabilizer}, the stabilizer
$H\le\Gamma$ of $z_0$ is cyclic. A generator $h$ fixes a label by the
hypothesis. Every power of $h$ fixes that same label, so $F^H$ is nonempty.
Let $y_0$ be its first member in the fixed finite ordering.

For a class $z=g\cdot z_0$, define $\ell(z)=g\cdot y_0$. This is
well-defined. If also $z=g'\cdot z_0$, then $g'^{-1}g\in H$, and hence
$g\cdot y_0=g'\cdot y_0$. The formula is visibly equivariant.

Finally, for $\vec q\in\ttuple{Q_\lambda}{m}$, take any representatives
$\vec a$ and set
\[
                          L(\vec q)=\ell([w_{\vec a}]_{\mathrm{st}}).
\]
The preceding discussion shows independence of every representative and
coordinate equivariance. The construction gives a unique value by a formula
using only the stated parameters, so Separation and Replacement produce the
required set-function. Choice of representatives was used only to verify
that the formula has a value; no representative enters the definition.
\end{proof}

\begin{theorem}[Prikry classification, including finite families]\label{thm:classification}
In the Prikry setup, the three assertions in Theorem~\ref{thm:main}(2) are
equivalent.
\end{theorem}

\begin{proof}
The implication from a finite family to the fixed-point condition is
Theorem~\ref{thm:necessity}. The condition yields a single $\OD_{V_\kappa}$
rule by Proposition~\ref{prop:sufficiency}. Its low-rank parameters are all
in $W$ because $V_\kappa^V=V_\kappa^W$. The singleton containing this rule
is therefore a finite $\OD_{p,q_*}$ family for a suitable combined ground
parameter $p$.
\end{proof}

The phrase ``for some $p\in W$'' is important in the sufficiency statement.
For a \emph{prescribed} parameter $p$, no well-order of the reals is assumed
to be definable from $p,q_*$. Necessity holds for every prescribed ground
parameter. Sufficiency may adjoin the one explicitly stated low-rank
well-order parameter.

\begin{remark}[The five-label example]
Let $\Gamma=S_3$, and let $F$ be the disjoint union of the natural three-point
action and the two-point sign action. The identity fixes every label. A
transposition fixes the remaining point of the natural action. A three-cycle
acts trivially on the sign action, so fixes both of its labels. Thus every
element has a fixed label. No label is fixed by the entire group: the
three-point action has no common fixed point, and a transposition exchanges
the two sign labels. The classification nevertheless supplies a rule.
This example, already identified in the archive, remains a useful check
against an incorrectly strengthened necessity argument.
\end{remark}

\section{Selectors, tournaments, and linear orders}\label{sec:selectors}

For $0<r<n$, set
\[
 \Sel_{n,r}(Q)=\{f:[Q]^n\to[Q]^r:
                         (\forall B\in[Q]^n)\ f(B)\subseteq B\}.
\]
These sets are nonempty in the ambient $\ZFC$ universe: any well-order of
$Q$ provides such selectors. Our obstruction concerns definable
\emph{families} of them, not their ordinary existence.

\begin{theorem}[No finite definable family of proper selectors]\label{thm:selectors}
For every ground parameter $p\in W$, every $n\ge2$, and every $0<r<n$,
there is no nonempty finite family
$\mathscr F\subseteq\Sel_{n,r}(Q_\kappa)$ in $\OD_{p,q_*}$.
\end{theorem}

\begin{proof}
Each selector $f$ canonically determines a rule on ordered distinct tuples:
\begin{equation}\label{eq:selectorlabel}
 L_f(x_0,\ldots,x_{n-1})
       =\{i<n:x_i\in f(\{x_0,\ldots,x_{n-1}\})\}\in[n]^r.
\end{equation}
With $\Gamma=S_n$ acting naturally on $F=[n]^r$, this rule is equivariant.
Indeed, reindexing the same unordered input by a permutation changes the
indices of the selected points by exactly that permutation.

If $\mathscr F$ were a nonempty finite $\OD_{p,q_*}$ family, its image under
the explicit map $f\mapsto L_f$ would be a nonempty finite definable family
of equivariant rules. Theorem~\ref{thm:necessity} would say that every
permutation of $n$ fixes an $r$-element subset of $n$. But an $n$-cycle is
transitive: any nonempty invariant subset contains the entire orbit of
any one of its elements, and hence is all of $n$. No proper nonempty
$r$-element subset is invariant. This contradicts $0<r<n$.
\end{proof}

\subsection{A direct view of the amplification}

For this particular theorem one can see the last part of the proof without
ordered labels. If $|\mathscr F|=M$, put $H=\lcm(1,\ldots,M)$ and take a
single cycle of $nH$ quotient classes. Encode each selector by its table on
all $n$-subsets of these classes. After deciding the finite set of tables,
deleting one late generic point rotates that set of tables. Its $H$th
power fixes every table. On the block
\[
                         B_0=\{0,H,2H,\ldots,(n-1)H\},
\]
the input rotation by $H$ is an $n$-cycle. A fixed table must therefore
select a subset of $B_0$ invariant under this transitive cycle, so it selects
either none or all of $B_0$. This is incompatible with a proper $r$-selection.
The larger configuration is what prevents a finite family from evading the
argument by permuting its members.

\begin{corollary}[Order-theoretic consequences]\label{cor:orders}
There is no nonempty finite $\OD_{p,q_*}$ family of tournaments on
$Q_\kappa$, of linear orders on $Q_\kappa$, or of choice functions on its
nonempty finite subsets. This holds for every $p\in W$.
\end{corollary}

\begin{proof}
A tournament canonically selects the winner of each two-element subset.
A linear order canonically selects its smaller element. A finite-set
choice function restricts to a selector on two-element subsets. Each of
these is a definable map from the relevant class of objects to
$\Sel_{2,1}(Q_\kappa)$. The image of a nonempty finite definable family under
such a map is again a nonempty finite definable family, contradicting
Theorem~\ref{thm:selectors}. It is immaterial if several original objects
have the same image.
\end{proof}

In particular these obstructions hold in $\OD_{V_\kappa}$ and, more
strongly, in $\OD_{V_\kappa,q_*}$. The additional parameter $q_*$ does not
supply a privileged representative or a starting index for its tail.

\section{Relative topology and full phase saturation}\label{sec:topology}

The same recoding lemma has a local consequence that is stronger than
cardinality alone. It identifies an actual neighborhood contained in every
definable trace through each of its points.

\subsection{Clopen definable traces}

For $P=(s,A)\in\PP_U^W$, define the basic subset of $q_*$
\[
      \mathcal N(s,A)=\{d\in q_*:(s,A)\in G_d\}.
\]
These sets form a basis: every $d\in q_*$ realizes conditions, and two
conditions realized by the same $d$ have a common stronger condition
realized by $d$. In detail, use the longer of their stems and intersect
the two measure-one sets above that stem. We call the resulting topology
on $q_*$ the \emph{relative Prikry topology}. Its definition uses the ground
forcing; it is not the product topology on enumerations of $q_*$.

\begin{theorem}[Local constancy]\label{thm:clopen}
Let $p\in W$.
Every $T\subseteq q_*$ in $\OD_{p,q_*}$ is clopen in the relative Prikry
topology. More generally, if $E\in W$ and
$h:q_*\to E$ is $\OD_{p,q_*}$, then $h$ is locally constant when $E$ has the
discrete topology.
\end{theorem}

\begin{proof}
Fix $b\in q_*$. By Lemma~\ref{lem:recoding}, $W[b]=V$ and $b$ is generic.
Suppose first that $b\in T$. In the forcing language, use the formula
that defines $T$ from $p,[\dot c]$ and the fixed ordinal parameters. The
truth lemma supplies a condition $P\in G_b$ forcing that its generic
sequence belongs to the uniquely defined set. For any $d\in\mathcal N(P)$,
we have $d\in q_*$ and $W[d]=V$. The formula therefore defines the original
set $T$, and $P\in G_d$ forces $d\in T$. Thus
$\mathcal N(P)\subseteq T$.

If $b\notin T$, apply the identical argument to $q_*\setminus T$.
This proves openness of both $T$ and its complement.

For the last assertion, $e=h(b)$ belongs to the ground-model set $E$,
so $e\in W$. Apply the truth lemma to the statement that the uniquely
defined function takes the value $\check e$ at the generic sequence.
The same argument gives a neighborhood $\mathcal N(P)$ of $b$ on which
$h$ is constantly $e$. No assumption that $E$ is finite or small is needed.
\end{proof}

An ordinal-valued definable map on $q_*$ has the same property. One can
apply the last statement with a ground ordinal larger than its range,
since the ground and the extension have the same ordinals.

\subsection{Every nonempty neighborhood has every phase}

For $a,b\in q_*$ define the integer phase
\begin{equation}\label{eq:phase}
             \ch_a(b)=|b\setminus a|-|a\setminus b|\in\mathbb Z.
\end{equation}
Adding one new point to $b$ raises its phase by one; deleting one point
lowers it by one. These identities hold even if the edited point belongs
to the reference set $a$. They follow either by separating the two finite
symmetric differences or by comparing the lengths of prefixes preceding
a common tail. Also
\[
        \ch_a(b)+\ch_b(d)=\ch_a(d)
                      \qquad(a,b,d\in q_*).
\]

\begin{lemma}[Phase-coordinate geometry of a basic neighborhood]
\label{lem:geometry}
Suppose $\mathcal N(s,A)$ is nonempty, put $\ell=|s|$, and fix
$a\in q_*$. For every $r\in\mathbb Z$ and every $n\ge\ell$,
\[
 \{b(n):b\in\mathcal N(s,A),\ \ch_a(b)=r\}
\]
contains a tail of $A$. In particular its cardinality is $\kappa$.
Moreover $|q_*|=\kappa$.
\end{lemma}

\begin{proof}
Fix $n\ge\ell$. Extend the finite sequence $s$ by $n-\ell$ entries from $A$,
obtaining a sequence $s'$ of length $n$. When $n=0$, use the empty
sequence and any harmless lower bound in place of its last point. For
each $\xi\in A$ above all entries of $s'$, let
$v=s'{}^\frown\seq{\xi}$, a sequence of length $n+1$.

Choose a sufficiently late tail of $a$ whose points all lie in $A$ and
all exceed $\xi$; this is possible because $a$ is Prikry generic. Append
this tail to $v$, obtaining $b_0\in\mathcal N(s,A)\cap q_*$ with
$b_0(n)=\xi$. Let $r_0=\ch_a(b_0)$.

If $r\ge r_0$, insert $r-r_0$ fresh points of $A\setminus b_0$ above $\xi$.
There are enough: $A$ has cardinality $\kappa$, whereas $b_0$ is countable,
and bounded initial segments have size less than $\kappa$. If $r<r_0$,
delete $r_0-r$ points of the infinite tail of $b_0$ above $\xi$. In either
case the resulting $b$ still has initial segment $v$, still realizes
$(s,A)$, and differs only finitely from $a$. The phase identities give
$\ch_a(b)=r$, while $b(n)=\xi$ is unchanged.

This realizes every $\xi$ in the indicated tail of $A$. Ground-model
measure-one sets have ground size $\kappa$, and cardinal preservation
keeps that size in $V$. Thus the projection has size $\kappa$.

Finally, every representative of $q_*$ is a finite symmetric change of
$a$, and every such change is still a cofinal set of order type $\omega$.
The bijection $F\mapsto a\triangle F$ from $[\kappa]^{<\omega}$ to $q_*$
shows that $|q_*|=\kappa$.
\end{proof}

\begin{theorem}[Full phase and coordinate saturation]\label{thm:phase}
Let $T\subseteq D_\kappa$ be $\OD_{p,q_*}$ for $p\in W$, and suppose
$T\cap q_*\ne\varnothing$. For every $a\in q_*$ and every $r\in\mathbb Z$,
\[
                    |\{b\in T\cap q_*:\ch_a(b)=r\}|=\kappa.
\]
There is $n_0<\omega$ such that, simultaneously for every integer $r$ and
every $n\ge n_0$,
\[
          |\{b(n):b\in T\cap q_*,\ \ch_a(b)=r\}|=\kappa.
\]
The integer $n_0$ can be chosen independently of the reference $a\in q_*$.
\end{theorem}

\begin{proof}
By Theorem~\ref{thm:clopen}, the nonempty trace contains a nonempty basic
neighborhood $\mathcal N(s,A)$. Put $n_0=|s|$. Lemma~\ref{lem:geometry}
gives all of the lower bounds, for any chosen $a$ and any integer phase.
The cardinality $|q_*|=\kappa$ and the fact that coordinates lie below
$\kappa$ give the upper bounds. The same neighborhood works for every $a$.
\end{proof}

Notice that cofinality alone would not give the asserted cardinality of a
projection: $\cf^V(\kappa)=\omega$. The proof instead exhibits a tail of a
ground-model measure-one set in each projection. This is why the conclusion
is full size $\kappa$, not merely unboundedness.

\section{Rigidity and the one-measurable upper bound}\label{sec:consistency}

For completeness, we recover the rigidity of $q_*$ by the cone argument
already used in the archive. This also verifies that the model used for
the new finite-symmetry conclusions is the same kind of model as its
one-measurable rigidity construction.

\subsection{Tail-class-preserving cone isomorphisms}

\begin{lemma}[No new definable sets of ordinals from the tail class]
\label{lem:groundordinals}
If $p\in W$ and $x$ is a set of ordinals in $\OD_{p,q_*}^V$, then $x\in W$.
\end{lemma}

\begin{proof}
Given conditions $(s,A)$ and $(t,B)$, choose
$C\in U$ with $C\subseteq A\cap B$ and all its points above both stems.
The cones below $(s,C)$ and $(t,C)$ are isomorphic by replacing the stem
$s$ with $t$ and leaving every subsequently appended finite sequence and
measure-one tail unchanged. Thus
\[
                    (s{}^\frown u,D)\longmapsto(t{}^\frown u,D).
\]
The generic sequences in corresponding extensions differ by the finite
stems $s,t$. Consequently the cone isomorphism preserves all check names,
including $p$ and ordinals, and preserves the name for the \emph{tail class}
of the generic sequence. This last assertion also follows directly by
interpreting the two names in corresponding generic extensions: they
produce the same universe and the same equivalence class.

It follows that the top condition decides every sentence of the form
$\psi(\check p,[\dot c],\check{\vec\alpha})$. Otherwise two conditions
would force opposite values. Pass to the isomorphic sub-cones just
described. The isomorphism transfers the first forced value to the
second cone, contradicting its opposite forced value.

Choose an ordinal $\theta$ with $x\subseteq\theta$, and fix a formula
defining $x$ uniquely from $p,q_*,\vec\alpha$. For each $\beta<\theta$,
apply the preceding conclusion to the assertion that the uniquely defined
object exists and contains $\beta$. In $W$, Separation using the forcing
relation forms the set of $\beta<\theta$ for which the top condition
forces that assertion. In the actual extension this set is precisely
$x$, by uniqueness and the forcing theorem. Hence $x\in W$.
\end{proof}

\subsection{Short cofinal families cannot be definable}

Let
\[
 \mathcal C_\kappa=
   \{a\subseteq\kappa:\sup a=\kappa,\ |a|^V<\kappa\}.
\]
Recall that rigidity of $q_*$ means that every nonempty
$\OD_{V_\kappa,q_*}$ subfamily of $\mathcal C_\kappa$ has size at least
$\kappa$.

\begin{proposition}[Rigidity, with ground parameters]\label{prop:rigidity}
For every $p\in W$, every nonempty
$\mathscr A\subseteq\mathcal C_\kappa$ in $\OD_{p,q_*}$ has
$|\mathscr A|^V\ge\kappa$. In particular $q_*$ is rigid.
\end{proposition}

\begin{proof}
Suppose instead that $0<|\mathscr A|^V<\kappa$. Let $\eta$ be the least
order type of a member of $\mathscr A$, and put
\[
             \mathscr A_0=\{a\in\mathscr A:\ot(a)=\eta\},
             \qquad u=\bigcup\mathscr A_0.
\]
Every member of $\mathcal C_\kappa$ has order type below the initial ordinal
$\kappa$, so $\eta<\kappa$. The set $u$ is unbounded in $\kappa$, is
$\OD_{p,q_*}$, and satisfies
\[
                      |u|^V\le |\mathscr A|^V\cdot|\eta|^V<\kappa.
\]
The last inequality holds for the product of two cardinals below an
infinite cardinal, even though $\kappa$ is singular in $V$.

By Lemma~\ref{lem:groundordinals}, $u\in W$. Unboundedness of a given set of
ordinals in $\kappa$ is absolute between these two transitive models.
Since $\kappa$ is regular in $W$, an unbounded ground set has ground size
$\kappa$. A ground bijection witnessing that size remains a bijection in
$V$, contradicting the displayed inequality.

Every finite tuple of parameters from $V_\kappa^V$ belongs to $W$, so the
last assertion follows.
\end{proof}

\begin{corollary}[Regularity in the relativized inner model]\label{cor:regular}
The cardinal $\kappa$ is regular in
$N_{q_*}=\HOD_{V_\kappa\cup\{q_*\}}$, although $\cf^V(\kappa)=\omega$.
\end{corollary}

\begin{proof}
If $N_{q_*}$ had a cofinal map $f:\delta\to\kappa$ for some $\delta<\kappa$,
then $f$ and its range would be ordinal-definable from $q_*$ and finitely
many individual parameters in $V_\kappa$. The range would be an unbounded
set of size at most $|\delta|<\kappa$. Lemma~\ref{lem:groundordinals} places
it in $W$, giving the same contradiction to ground regularity and cardinal
preservation as above. This proof does not assume Choice in $N_{q_*}$.
\end{proof}

\subsection{What consistency conclusion has actually been proved}

\begin{corollary}[One-measurable construction]\label{cor:consistency}
If $\ZFC$ with a measurable cardinal is consistent, then so is $\ZFC$ with
a strong limit cardinal $\kappa$ of cofinality $\omega$ and a rigid
$q_*\in Q_\kappa$ satisfying all conclusions of Theorem~\ref{thm:main}.
\end{corollary}

\begin{proof}
A measurable cardinal carries a normal measure; force with its normal
Prikry forcing. Input~\ref{in:prikry} gives the strong-limit, cofinality,
and rank-agreement assertions. Proposition~\ref{prop:rigidity} gives
rigidity. Theorems~\ref{thm:classification}, \ref{thm:selectors}, and
\ref{thm:phase} give the finite-symmetry and saturation conclusions.
The ultrafilter assertion is established in Section~\ref{sec:ultrafilters}.
These proofs are forcing arguments over $\ZFC$; the transitive-model
notation used for exposition imposes no additional transitive-model
consistency assumption.
\end{proof}

The supplied synthesis gives an exact measurable calibration for a weaker
package including rigidity, using quoted core-model covering theorems.
The construction here extends its \emph{upper-bound} side to the finite
selector and finite-label obstructions. No new lower-bound proof is offered
here, and none is needed for the new same-universe recoding theorem.
In particular, no part of the argument asserts that a model with an
ultraexacting cardinal is obtainable from one measurable cardinal.

\section{A sharp countable threshold for local ultrafilters}\label{sec:ultrafilters}

An ultrafilter in this section is an ultrafilter on the full Boolean algebra
$\Pow^V(q)$ of subsets of a tail class $q$. This is not a normal measure on
$\kappa$ inside an inner model. Keeping these two domains separate is
essential.

\subsection{No finite definable family at the Prikry class}

\begin{theorem}[Finite local ultrafilter obstruction]\label{thm:nofiniteuf}
For every $p\in W$, there is no nonempty finite $\OD_{p,q_*}$ family of
ultrafilters on $\Pow^V(q_*)$.
\end{theorem}

\begin{proof}
Suppose $\mathscr U$ is such a family, with $M=|\mathscr U|$. Fix an integer
$N>M$. For a representative $a\in q_*$, partition $q_*$ into the phase
residue classes
\[
             P_i(a)=\{b\in q_*:\ch_a(b)\equiv i\pmod N\},
                         \qquad i\in\mathbb Z/N\mathbb Z.
\]
Every ultrafilter contains exactly one cell of a finite partition. Thus
\[
 S(a)=\{i\in\mathbb Z/N\mathbb Z:
                  P_i(a)\in\mathcal U\text{ for some }\mathcal U\in\mathscr U\}
\]
is nonempty and has size at most $M$.

Apply Lemma~\ref{lem:decision} to the finite observation $S(a)$. Let
$P\in G_c$ decide its value to be $S(c)$, and delete one point of $c$
after the stem of $P$, obtaining $d$. The recoding lemma leaves $V,q_*$,
and the definable family $\mathscr U$ unchanged, and yields $S(d)=S(c)$.
But $\ch_c(d)=-1$, so the phase cocycle gives
\[
                         \ch_d(b)=\ch_c(b)+1
                                  \qquad(b\in q_*).
\]
Therefore $P_i(d)=P_{i-1}(c)$ and $S(d)=S(c)+1$ modulo $N$. A nonempty
subset of $\mathbb Z/N\mathbb Z$ invariant under addition of $1$ is the
whole set. Hence $|S(c)|=N>M$, a contradiction.
\end{proof}

\begin{corollary}[Finite global-kernel obstruction]\label{cor:nofinitekernels}
There is no nonempty finite $\OD_{p,q_*}$ family of functions $K$ on
$Q_\kappa$ such that $K(q)$ is an ultrafilter on $\Pow^V(q)$ for every
$q\in Q_\kappa$.
\end{corollary}

\begin{proof}
Evaluate every member of the family at $q_*$. The resulting image is a
nonempty finite $\OD_{p,q_*}$ family of ultrafilters on $q_*$, contrary to
Theorem~\ref{thm:nofiniteuf}.
\end{proof}

\subsection{A countable local family in ordinary $\ZFC$}

The positive construction does not require Prikry forcing or a measurable
cardinal. It uses a free ultrafilter on the integers and all of its shifts,
so that a finite change of representative only reindexes the countable
family.

Fix a free ultrafilter $D$ on $\mathbb Z$ containing $\omega$. Such an
ultrafilter exists in $\ZFC$, by transporting a free ultrafilter on $\omega$
to $\mathbb Z$. For $k\in\mathbb Z$, let
\[
      D_k=(\sh_k)_*D,\qquad \sh_k(n)=n+k,
\]
where the pushforward of an ultrafilter under $f:X\to Y$ is
\[
                         f_*D=\{A\subseteq Y:f^{-1}[A]\in D\}.
\]
Every $D_k$ is free and contains $\omega$, since translates of $\omega$
differ from it by only finitely many integers.

For $a\in D_\lambda$, define a map $t_a:\mathbb Z\to[a]$ by
\[
 t_a(n)=\begin{cases}
           \{a(j):j\ge n\},&n\ge0,\\
           a,&n<0.
         \end{cases}
\]
Its values at negative indices are irrelevant to all the $D_k$. Set
\begin{equation}\label{eq:countuf}
                  \mathscr U_D([a])=
                         \{(t_a)_*D_k:k\in\mathbb Z\}.
\end{equation}

\begin{theorem}[Countable tail-supported ultrafilter correspondence]
\label{thm:countuf}
For every uncountable cardinal $\lambda$ of cofinality $\omega$, formula
\eqref{eq:countuf} defines a function on $Q_\lambda$, independent of the
chosen representative $a$. For every $q\in Q_\lambda$, the set
$\mathscr U_D(q)$ is countably infinite and consists of free ultrafilters
on the full $\Pow(q)$. Every one of these ultrafilters contains the tail
ray
\[
                       R_b=\{t_b(n):n<\omega\}
                       \qquad\text{for every }b\in q.
\]
The function is definable from $D$ and $\lambda$, and its value at $q$ is
$\OD_{D,q}$. The parameter $D$ belongs to $V_\lambda$.
\end{theorem}

\begin{proof}
\textbf{Existence and full domain.}
For fixed $a,k$, the pushforward $(t_a)_*D_k$ is an ultrafilter on the
entire $\Pow([a])$: for every subset $B\subseteq[a]$, its inverse image is
a subset of $\mathbb Z$, and $D_k$ decides that inverse image or its
complement. This is not an ultrafilter merely on a specified subalgebra.
The map $t_a$ is injective on $\omega$, so the inverse image of a singleton
meets $\omega$ in at most one point. Because $D_k$ is free and contains
$\omega$, no singleton belongs to the pushforward.

\textbf{Independence of the representative.}
If $b=^*a$, their increasing enumerations have synchronized tails. More
explicitly, choose $I,J<\omega$ such that
\[
                            a(I+n)=b(J+n)
                                    \qquad(n<\omega).
\]
With $e=I-J$, this gives $t_b(n)=t_a(n+e)$ for all sufficiently large
$n$. Free ultrafilters containing $\omega$ disregard the exceptional
indices. Therefore
\[
                           (t_b)_*D_k=(t_a)_*D_{k+e}
                                   \qquad(k\in\mathbb Z).
\]
Since $k\mapsto k+e$ permutes $\mathbb Z$, the sets of pushforwards in
\eqref{eq:countuf} are equal. This proves independence and also shows
that the value is definable from $D,q$, without an individual
representative as a parameter.

\textbf{Distinctness of all shifts.}
No free ultrafilter on $\mathbb Z$ is fixed by a nonzero integer
translation. To see this, suppose it were fixed by translation by
$h>0$. Let
\[
          E_h=\{n\in\mathbb Z:n\bmod 2h\in\{0,1,\ldots,h-1\}\}.
\]
Translation by $h$ exchanges $E_h$ and its complement. Invariance would
make an ultrafilter contain both or neither, which is impossible.
A negative period implies the corresponding positive period. It follows
that $D_k\ne D_l$ whenever $k\ne l$.

Pushforward by $t_a$ preserves this distinction. Indeed, all $D_k$
concentrate on $\omega$, where $t_a$ is injective. For $B\subseteq\omega$,
the inverse image of $t_a[B]$ agrees with $B$ on $\omega$. Consequently
one can recover every membership decision of $D_k$ on $\omega$ from its
pushforward, and decisions on $\omega$ determine an ultrafilter containing
$\omega$. Thus the pushforwards are pairwise distinct. Formula
\eqref{eq:countuf} gives exactly $\aleph_0$ of them.

\textbf{Tail support.}
Each $(t_a)_*D_k$ contains $R_a$, since $t_a^{-1}[R_a]$ contains $\omega$.
If $b=^*a$, the synchronized-tail identity shows that $R_a$ and $R_b$
differ by finitely many elements of $q$. A free ultrafilter containing
$R_a$ therefore contains $R_b$ as well.

\textbf{Rank and definability.}
An ultrafilter on $\mathbb Z$ is a set of subsets of a countable
low-rank set. Its rank, under the usual coding of integers, is below
$\omega+\omega$. Hence $D\in V_\lambda$. Independence of representatives
makes \eqref{eq:countuf} a functional definition on the set $Q_\lambda$;
Replacement supplies its graph. All the claimed definability statements
now follow from the explicit formula.
\end{proof}

\begin{corollary}[Exact local threshold]\label{cor:threshold}
In the Prikry setup, the least size of a nonempty
$\OD_{V_\kappa,q_*}$ family of ultrafilters on $\Pow(q_*)$ is $\aleph_0$.
There is such a countable family all of whose ultrafilters are free and
contain every representative's tail ray.
\end{corollary}

\begin{proof}
Theorem~\ref{thm:nofiniteuf} rules out every positive finite size, since
all permitted low-rank parameters belong to $W$. Choose $D$ as in
Theorem~\ref{thm:countuf}. It belongs to $V_\kappa$, and that theorem gives
a countably infinite family with all the required properties.
\end{proof}

\subsection{Why this does not solve the global countable-kernel question}

Theorem~\ref{thm:countuf} gives one definable \emph{correspondence}
\[
       q\longmapsto\mathscr U_D(q),
              \qquad |\mathscr U_D(q)|=\aleph_0.
\]
It does not give a definable countable set of functions
$\{K_n:n<\omega\}$ with $K_n(q)$ an ultrafilter on $q$ simultaneously
for all $q$. The displayed indexing in \eqref{eq:countuf} uses a
representative $a$. Replacing $a$ changes that indexing by an integer
shift which can depend on $q$. The \emph{set} of values is canonical,
but the index origin is not.

Ordinary Choice can select representatives and synchronize these
indexings, but it does not make the resulting synchronization definable
from $D$ or other permitted parameters. Therefore the countably infinite
global-family question in the supplied synthesis is not resolved by
this construction. The local threshold in Corollary~\ref{cor:threshold}
is an exact theorem of a different, explicitly stated kind.

\section{Proof audit and the remaining boundary}\label{sec:audit}

\subsection{Dependencies and new content}

The new forcing argument depends on normal Prikry preservation, the
Mathias criterion, and the ordinary forcing theorem. It does not depend
on an ultraexacting witness, a correctness-height argument, an iterated
ultrapower, a core-model covering theorem, or the strongly compact
sandwich results in the supplied reports.

The combinatorial amplification and the incidence-word sufficiency
construction have antecedents in the archive. The new step is the
finite-table recoding lemma in a normal Prikry extension, together with
its applications to finite families of arbitrary finite equivariant
rules. The local clopen theorem and its full-sized phase-coordinate
corollary isolate a further consequence of the same mechanism. The
countable ultrafilter correspondence supplies a sharp local positive
counterpart to the finite obstruction.

This is new relative to the mathematical development in the supplied
archive: it proves the previously unproved Prikry-model assertion asked
there and strengthens it to allow both the generic quotient parameter
and arbitrary fixed ground-model parameters. The literature check
verified the classical forcing inputs and the large-cardinal comparison
used for context. It did not establish bibliographic priority for every
application or rule out equivalent formulations elsewhere.

\subsection{Checks that prevent false strengthenings}

\begin{assessment}{Six distinctions maintained throughout the proofs}
\textbf{One class versus one representative.}
The parameter $q_*$ is fixed under finite recoding; $c$ is not available
as a defining parameter. Allowing $c$ would destroy the central symmetry.

\textbf{A definable family versus definable members.}
The finite set of evaluation tables is decided as a set. Members of the
original family need not be individually definable or fixed.

\textbf{A Prikry model versus abstract rigidity.}
We prove the finite obstructions in a concrete forcing extension and
also prove that its generic class is rigid. We do not deduce the
recoding property from rigidity alone.

\textbf{Elementwise versus common fixed points.}
Each $g\in\Gamma$ must fix some label, which may depend on $g$. The
five-label $S_3$ example prevents replacing this with a common fixed
label requirement.

\textbf{Local ultrafilters versus inner-model measures.}
The ultrafilters in Section~\ref{sec:ultrafilters} decide all subsets
of $q_*$ in $V$. They are not normal measures on $\kappa$ in a
relativized $\HOD$.

\textbf{Countably many values versus countably many global functions.}
A countable-valued correspondence on $Q_\lambda$ does not provide a
definable simultaneous enumeration of its fibres.
\end{assessment}

\subsection{Questions that are still genuinely open in this continuation}

The next structural problem is whether the finite-table recoding
conclusion follows from some intrinsic strengthening of rigidity,
weaker than a full normal Prikry presentation. This would distinguish
which part of the phenomenon is topological and which part is forced
by regularity in $N_q$.

The unrestricted implication from rigidity to the finite-selector
obstruction remains unproved here. So do the questions about ordinary
exactingness, countably infinite definable families of global selectors,
and countably infinite definable families of global ultrafilter-valued
kernels at an ultraexacting cardinal. Neither the positive local
ultrafilter theorem nor the one-measurable construction decides them.

The proved advance is therefore specific but substantial: the source
archive's finite-symmetry boundary is no longer confined to the
ultraexacting construction. It is already forced by finite recodings
of an ordinary normal Prikry generic sequence.

\appendix
\section{Finite verification and reproducibility}\label{app:tests}

The archive's Lean file
\nolinkurl{Cardinals/Combinatorics/FiniteCycles.lean} separates finite
permutation amplification from the large-cardinal interface. That
separation guided Section~\ref{sec:tables}. No new Lean theorem is
claimed here, and the supplied Lean development was not recompiled for
this report.

The accompanying Python script \nolinkurl{finite_checks.py} uses only the
standard library. It performs exact finite checks of three ingredients:
it exhaustively computes cyclic orbits of small selector tables, checks
the amplified permutation embedding for every permutation of at most
six labels and family-size bounds at most five, and verifies the
five-label $S_3$ fixed-point example. It also checks that finite-word
stabilizers with faithful label separation are cyclic, for the small
periodic examples described in its output.

The checks are deterministic. They neither simulate a measurable
cardinal nor verify a forcing argument, and they are not a substitute
for the English proofs. Their purpose is to detect an incorrect
permutation convention, an insufficient amplification length, or an
accidental common-fixed-point assumption.

\begin{center}
\small
\renewcommand{\arraystretch}{1.13}
\begin{tabularx}{\textwidth}{@{}P{.29\textwidth}P{.22\textwidth}Y@{}}
\toprule
\textbf{Finite test}&\textbf{Exact value}&\textbf{Interpretation}\\\midrule
Binary selectors on 2 points&2 tables; minimum cyclic orbit 2&The two
choices can be exchanged; the unamplified test does not exclude a
2-member family\\
Binary selectors on 4 points&64 tables; minimum cyclic orbit 4&The
4-cycle amplification excludes invariant families of size at most 2\\
Binary selectors on 6 points&32768 tables; minimum cyclic orbit 2&An
arbitrary large cycle is insufficient; its length must match the
amplification construction\\
Three-to-one selectors on 3 points&3 tables; minimum cyclic orbit 3&A
3-cycle has no invariant singleton choice\\
Five-label $S_3$ target&Each of 6 permutations fixes a label; no common
fixed label&The theorem uses the elementwise condition, not a global
fixed point\\
\bottomrule
\end{tabularx}
\end{center}

The distribution of all orbit sizes and all other test counts are
recorded in \nolinkurl{finite_checks.txt}. The script's assertions fail
with a nonzero exit status if a check disagrees with the expected
identity. Its input sizes are deliberately bounded so that it runs
quickly on an ordinary Python installation.

\subsection*{Build}

The report uses the same \texttt{newpxtext}/\texttt{newpxmath} font setup,
page geometry, theorem style, and Forest--Olive--Sage color definitions
as the supplied synthesis. A standard TeX installation with those
packages can rebuild the PDF by running
\begin{verbatim}
pdflatex -interaction=nonstopmode -halt-on-error Prikry_Finite_Symmetry.tex
pdflatex -interaction=nonstopmode -halt-on-error Prikry_Finite_Symmetry.tex
python3 finite_checks.py
\end{verbatim}
An additional LaTeX pass may be used if a local engine requests it for
updated cross-references. Font files are not redistributed.

\begin{thebibliography}{9}

\bibitem{archive}
\emph{Large cardinals at the inconsistency frontier: a synthesis}.
Supplied research archive \nolinkurl{Cardinals4.zip},
\nolinkurl{docs/Large_Cardinals_Synthesis.tex}, dated 18 September 2026.
Especially Sections~9, 12.3, 12.5, 13.4--13.5, and the final question of
Section~15. Consulted together with
\nolinkurl{docs/Large_Cardinals_Unified_Report.tex} and the accompanying
\nolinkurl{Cardinals/} Lean sources. Internal research material; not
presented here as an independently peer-reviewed source.

\bibitem{mathias}
A.~R.~D.~Mathias,
\emph{On sequences generic in the sense of Prikry},
Journal of the Australian Mathematical Society \textbf{15} (1973), 409--414.
\doi{10.1017/S1446788700028755}.

\bibitem{benhamou}
Tom Benhamou,
\emph{Prikry forcing and tree Prikry forcing of various filters},
Archive for Mathematical Logic \textbf{58} (2019), 787--817.
\doi{10.1007/s00153-019-00660-3}.
Open version: \arxiv{1801.04424}, version~2, 25 September 2018.
The theorem numbers in this report refer to that open version.

\bibitem{abgl}
Juan P.~Aguilera, Joan Bagaria, Gabriel Goldberg, and Philipp L\"ucke,
\emph{Large cardinals beyond HOD},
\arxiv{2509.10254}, version~1, 12 September 2025.
Used for the contextual comparison with $\Izero$, not as an input to the
new finite-tail recoding proofs.

\bibitem{abl}
Juan P.~Aguilera, Joan Bagaria, and Philipp L\"ucke,
\emph{Large cardinals, structural reflection, and the HOD Conjecture},
\arxiv{2411.11568}, version~4, 16 September 2025.
Background for the exacting and ultraexacting framework continued by
the supplied archive.

\end{thebibliography}
\end{document}
